\documentclass[11pt]{article}

\usepackage[a4paper,margin=2.2cm]{geometry}
\usepackage[T1]{fontenc}
\usepackage[utf8]{inputenc}
\usepackage{lmodern}
\usepackage{amsmath,amssymb,mathtools,bm}
\usepackage{physics}
\usepackage{microtype}
\usepackage[dvipsnames]{xcolor}
\usepackage[most]{tcolorbox}
\usepackage{enumitem}
\usepackage{titlesec}

% ---------- Section formatting ----------
\titleformat{\section}{\Large\bfseries}{\thesection}{0.75em}{}
\titleformat{\subsection}{\large\bfseries}{\thesubsection}{0.75em}{}
\titlespacing*{\section}{0pt}{1.2ex plus .2ex minus .2ex}{0.8ex}
\titlespacing*{\subsection}{0pt}{1.0ex plus .2ex minus .2ex}{0.5ex}

% ---------- Step box environment ----------
\newcounter{derivstep}

% ---------- Assumption box environment ----------
\newcounter{assumption}

\newtcolorbox{assumptionbox}[1]{
    breakable,
    enhanced,
    colback=green!2,
    colframe=green!45!black,
    boxrule=0.7pt,
    arc=1.5mm,
    left=2mm,right=2mm,top=1mm,bottom=1mm,
    title={\bfseries Assumption \theassumption\;--\;#1},
    coltitle=black,
    colbacktitle=green!10,
    fonttitle=\bfseries
}

\newenvironment{assumptionenv}[1]
{
    \refstepcounter{assumption}
    \begin{assumptionbox}{#1}
}
{
    \end{assumptionbox}
}

\newtcolorbox{stepbox}[1]{
    breakable,
    enhanced,
    colback=blue!2,
    colframe=blue!55!black,
    boxrule=0.7pt,
    arc=1.5mm,
    left=2mm,right=2mm,top=1mm,bottom=1mm,
    title={\bfseries Step \thederivstep\;--\;#1},
    coltitle=black,
    colbacktitle=blue!10,
    fonttitle=\bfseries
}

\newenvironment{derivationstep}[1]
{
    \refstepcounter{derivstep}
    \begin{stepbox}{#1}
}
{
    \end{stepbox}
}

% ---------- Source note ----------
\newcommand{\sourcenote}[1]{%
    \par\vspace{0.4em}%
    \textcolor{gray}{\footnotesize\textbf{Source:} #1}%
}

% ---------- Shortcuts ----------
\newcommand{\R}{\mathbb{R}}
\newcommand{\Ac}{A_c}
\newcommand{\Bc}{B_c}
\newcommand{\Cc}{C_c}
\newcommand{\Dc}{D_c}
\newcommand{\Ad}{A_d}
\newcommand{\Bd}{B_d}
\newcommand{\Td}{T_d}

\begin{document}

\begin{center}
    {\LARGE \textbf{Formatted Derivation Notes}}\\[0.4em]
    {\large Dual-Gantry LPV Discrete-Time State-Space Derivation}
\end{center}

\vspace{0.8em}

\section*{Model Classification}

Tóth \cite{toth2010} defines CT-LPV-SS representations as models of an underlying physical
system $S$ where the state-space matrices depend on a scheduling variable $p$:
\[
    \dot{x}(t) = A_c(p(t))\,x(t) + B_c(p(t))\,u(t), \qquad
    y(t) = C_c(p(t))\,x(t) + D_c(p(t))\,u(t).
\]

The dual-gantry model fits this definition. The state vector is
\[
    x = \begin{bmatrix} X & \Theta & Y & \dot{X} & \dot{\Theta} & \dot{Y} \end{bmatrix}^\top \in \R^6,
\]
so the payload Y-position is the third component: $Y = x_3$.
The scheduling variable is $p = Y$, and the matrices $A_c(Y)$ and $B_c(Y)$ depend on it
through the mass matrix $M(Y)$ (see Derivation Step~1 and Step~2).

\textbf{$C_c$ and $D_c$ are constant.}
In the general LPV definition all four matrices may depend on $p$. For our model:
\[
    C_c = \begin{bmatrix} I_3 & 0_{3\times 3} \end{bmatrix}, \qquad D_c = 0_{3\times 3}.
\]
Neither depends on $Y$. The LPV structure is carried entirely by the state equation; the
output equation is parameter-invariant. As a consequence, $C_d = C_c$ and $D_d = D_c$
follow trivially from the discretization with no scheduling complication.

The classification requires specifying the nature of the scheduling dependence, for which
Tóth distinguishes two cases:

\begin{itemize}
    \item \textbf{Static dependence}: $p$ is an exogenous signal, independent of the system
    state. The scheduling variable is externally measured or prescribed.
    \item \textbf{Dynamic dependence}: $p = f(x)$, i.e., the scheduling variable is a
    function of the system state. This is referred to as \emph{quasi-LPV}.
\end{itemize}

\textbf{Our system is quasi-LPV with dynamic scheduling dependence.}
The scheduling variable $p = Y = x_3$ is the third component of the state vector, so it is
not an exogenous signal but evolves as part of the system dynamics. Both classifications are
valid CT-LPV representations in Tóth's framework; quasi-LPV is a subclass of LPV, not
outside it. The distinction matters at the discretization step: the ZOH complete method
is exact for static dependence and approximate for dynamic dependence (see quasi-LPV caveat
in the assumptions below).

\vspace{0.8em}

\section*{Assumptions (Tóth \cite{toth2010}, pp.~5--6)}

The complete method for LPV discrete-time state-space conversion requires exactly two
assumptions to hold.

\begin{assumptionenv}{ZOH setting}
\textit{``We are given a CT-LPV system S, with CT input signal \(u_c\), scheduling signal
\(p_c\), and output signal \(y_c\), where \(u_c\) and \(p_c\) are generated by an ideal
ZOH\footnote{Tóth \cite{toth2010}, footnote~1: \textit{``The ZOH device is a signal hold
instrument providing a CT signal which is constant till the device is commanded to change it
to a new value in a piecewise constant manner.''}}
device and \(y_c\) is sampled in a perfectly synchronized manner with \(T_d > 0\) as the
sampling period or discretization time-step.''}

\medskip
The assumption has two parts: (i) the control input \(u_c\) is ZOH-held, and (ii) the
scheduling signal \(p_c\) is ZOH-held. Both must be piecewise-constant on
\([kT_d,\,(k+1)T_d)\).

\medskip
\textbf{Part (i): control input \(u_c\)}
The gantry runs a digital control loop at sample rate \(f_s\) (PLTI rate,
\cite{etel_accuret} -- to be verified). The sample period is
\[
    T_d = \frac{1}{f_s} = \frac{1}{16\,000} = 62.5\,\mu\text{s}.
\]
The force command \(u[k]\) is computed once per cycle by the discrete controller and held
constant until the next cycle. This is the definition of a ZOH digital actuator.
\textbf{Satisfied exactly by construction.} \checkmark

\medskip
\textbf{Part (ii): scheduling signal \(p_c = Y\)}
Here the assumption holds only approximately. \(Y = x_3\) is a continuous state, so it
evolves within the interval as the integrator runs. We treat it as ZOH-held by reading
\(Y[k] = x_3(kT_d)\) at the start of each sample. The maximum displacement of \(Y\)
within one interval is
\[
    \Delta Y_{\max} = v_{\max} \cdot T_d
    = 2\,\tfrac{\text{m}}{\text{s}} \times 62.5 \times 10^{-6}\,\text{s}
    = 1.25 \times 10^{-4}\,\text{m}
    = 0.125\,\text{mm},
\]
where \(v_{\max} = 2\,\text{m/s}\) is the maximum Y-axis velocity from \cite{etel_accuret},
confirmed in \texttt{main.m} (\texttt{vmax = 2}). The assumption therefore holds
approximately, with the intra-sample variation of \(Y\) treated as a frozen-at-sampling
approximation error. \checkmark\ (approximate; see quasi-LPV caveat below)
\end{assumptionenv}

\vspace{0.5em}

\begin{assumptionenv}{Switching effects}
\textit{``The switching behavior of the ZOH actuation has no effect on the CT plant, i.e.\
the switching of the signals is assumed to take place smoothly.''}

\medskip
Tóth notes: \textit{``this assumption is automatically satisfied in most numerical simulations
of LPV systems, like in the implemented numerical approaches of SIMULINK in MATLAB.''}

\medskip
The assumption requires that the step changes in \(u\) and \(p_c = Y\) at each sample
boundary are small enough not to excite the plant discontinuously.

\medskip
\textbf{Step size of \(u[k]\):}
The controller bandwidth is \(f_\text{bw} = 100\,\text{Hz}\) (\texttt{fbw = 100} in
\texttt{main.m}). The force command therefore changes on a timescale of
\[
    \tau_\text{ctrl} = \frac{1}{2\pi f_\text{bw}}
    = \frac{1}{2\pi \times 100}
    \approx 1.6\,\text{ms}
    = \frac{1.6 \times 10^{-3}}{T_d}
    \approx 25\;\text{samples}.
\]
Each sample-to-sample step in \(u\) is therefore approximately \(1/25\) of the full
control amplitude, which is a smooth, gradual change relative to the plant dynamics.

\medskip
\textbf{Step size of \(p_c = Y\):}
From Assumption~1, \(\Delta Y_{\max} = 0.125\,\text{mm}\) per sample. The operational
range of \(Y\) is \(700\,\text{mm}\) (800\,mm stroke with 5\% margin, \cite{etel_accuret}).
The per-sample step as a fraction of the full range is
\[
    \frac{\Delta Y_{\max}}{Y_{\text{range}}}
    = \frac{0.125}{700}
    \approx 1.8 \times 10^{-4}.
\]
The change in \(A_c(Y)\) per sample is therefore negligible: the matrix barely moves
between steps. \textbf{Satisfied.} \checkmark

\medskip
Our Python simulation uses numerical integration that mirrors the SIMULINK approach Tóth
explicitly endorses, providing an additional practical guarantee.
\end{assumptionenv}

\vspace{0.5em}

\begin{tcolorbox}[
    colback=orange!4, colframe=orange!60!black, boxrule=0.7pt, arc=1.5mm,
    left=2mm, right=2mm, top=1mm, bottom=1mm,
    title={\bfseries Quasi-LPV caveat}, coltitle=black, colbacktitle=orange!15,
    fonttitle=\bfseries
]
Tóth states the complete method is \emph{``only reasonable for the discretization of
LPV-SS representation with static dependence as dynamic dependence requires a higher-order
hold approach.''} This is not a formal assumption but an assessment of accuracy.

\medskip
Our scheduling variable \(Y = x_3\) is a system \emph{state} (dynamic dependence), so it is
not an exogenous signal but evolves as part of the system dynamics. The ZOH setting of
Assumption~1 is therefore applied approximately: \(Y\) is frozen at \(Y[k]\) rather than
truly held by an external ZOH device. This introduces a residual intra-sample error.

\medskip
\textbf{Timescale separation argument.}
The residual error is accepted as negligible based on the following argument.
The time for \(Y\) to traverse its full operational range \(Y_\text{range} = 700\,\text{mm}\)
at maximum speed is
\[
    t_\text{traverse} = \frac{Y_\text{range}}{v_{\max}}
    = \frac{0.7\,\text{m}}{2\,\text{m/s}}
    = 350\,\text{ms}
    = \frac{350 \times 10^{-3}}{T_d}
    = 5\,600\;\text{samples}.
\]
The plant's fastest relevant dynamics act on the closed-loop bandwidth timescale
\[
    \tau_\text{plant} = \frac{1}{2\pi f_\text{bw}}
    = \frac{1}{2\pi \times 100}
    \approx 1.6\,\text{ms}
    \approx 25\;\text{samples}.
\]
The timescale separation is therefore
\[
    \frac{t_\text{traverse}}{\tau_\text{plant}}
    = \frac{350\,\text{ms}}{1.6\,\text{ms}}
    \approx 220.
\]
Within any single plant timescale, \(Y\) moves at most
\(\tau_\text{plant} \times v_{\max} = 1.6\,\text{ms} \times 2\,\text{m/s} = 3.2\,\text{mm}\),
which is \(3.2/700 \approx 0.46\%\) of the full operational range. The frozen approximation
is therefore highly accurate in practice.

\medskip
\textbf{Numerical confirmation (Task~2.5).}
The relative change in \(A_c(Y)\) per sample is verified as
\(\|A_c(Y + \Delta Y_{\max}) - A_c(Y)\| / \|A_c(Y)\|\) at \(\Delta Y_{\max} = 0.125\,\text{mm}\).
\end{tcolorbox}

\vspace{1.2em}

\section*{Derivation}

\begin{derivationstep}{Equations of motion (logical coordinates)}

The gantry system with
\[
q = \begin{bmatrix} X & \Theta & Y \end{bmatrix}^{\top}
\]
and logical forces \(f_{\text{logical}}\) is given by
\[
M(Y)\ddot{q} + C_{\text{damp}}\dot{q} + Kq = f_{\text{logical}}.
\]

With
\[
M(Y)=
\begin{bmatrix}
m_1+m_2+m_b+m_h &
(m_1-m_2)\dfrac{L_b}{2} - m_h Y &
0
\\[0.6em]
(m_1-m_2)\dfrac{L_b}{2} - m_h Y &
J_b+J_h+(m_1+m_2)\dfrac{L_b^2}{4}+m_h d^2+m_h Y^2 &
-m_h d
\\[0.6em]
0 & -m_h d & m_h
\end{bmatrix},
\]

\[
C_{\text{damp}}=
\begin{bmatrix}
c_{g1}+c_{g2} &
(c_{g1}-c_{g2})\dfrac{L_b}{2} &
0
\\[0.6em]
(c_{g1}-c_{g2})\dfrac{L_b}{2} &
c_{b1}+c_{b2}+(c_{g1}+c_{g2})\dfrac{L_b^2}{4} &
0
\\[0.6em]
0 & 0 & c_y
\end{bmatrix},
\qquad
K=
\begin{bmatrix}
0 & 0 & 0 \\
0 & k_{b1}+k_{b2} & 0 \\
0 & 0 & 0
\end{bmatrix}.
\]

\sourcenote{\texttt{kamtin-fp-model/03 Simulink gantry/main.m} lines 12--54
(physical parameters and mass matrix).}

\end{derivationstep}

\begin{derivationstep}{First-order state-space (logical coordinates)}

Define the state
\[
x =
\begin{bmatrix}
q^{\top} & \dot{q}^{\top}
\end{bmatrix}^{\top}
=
\begin{bmatrix}
X & \Theta & Y & \dot{X} & \dot{\Theta} & \dot{Y}
\end{bmatrix}^{\top}
\in \R^6.
\]

Rewrite the second-order system as
\[
\ddot{q} = -M(Y)^{-1}C_{\text{damp}}\dot{q} - M(Y)^{-1}Kq + M(Y)^{-1}f_{\text{logical}}.
\]

Then
\[
\dot{x} = A_c(Y)x + B_c f_{\text{logical}},
\qquad
y = C_c x,
\]
with
\[
A_c(Y)=
\begin{bmatrix}
0_{3\times 3} & I_3 \\
-M(Y)^{-1}K & -M(Y)^{-1}C_{\text{damp}}
\end{bmatrix}
\in \R^{6\times 6},
\]
\[
B_c=
\begin{bmatrix}
0_{3\times 3} \\
M(Y)^{-1}
\end{bmatrix}
\in \R^{6\times 3},
\qquad
C_c=
\begin{bmatrix}
I_3 & 0_{3\times 3}
\end{bmatrix}
\in \R^{3\times 6},
\qquad
D_c = 0_{3\times 3}.
\]

\textbf{Note.}
\(A_c(Y)\) is singular. The top-left block is identically zero, and the zero rows of \(K\)
(the \(X\)- and \(Y\)-coordinates have no stiffness) propagate to \(\det(A_c)=0\).

\sourcenote{\texttt{kamtin-fp-model/03 Simulink gantry/functions/getss.m} (MATLAB ground truth).
Python implementations: \texttt{scripts/gantry/gantry\_ss.py} (numpy/scipy, validation only)
and \texttt{scripts/gantry/gantry\_lpv\_torch.py} (torch, differentiable, training loop).}

\end{derivationstep}

\begin{derivationstep}{Stage coordinate transform}

Stage forces and positions relate to logical coordinates via
\[
P=
\begin{bmatrix}
1 & 1 & 0 \\
\dfrac{L_b}{2} & -\dfrac{L_b}{2} & 0 \\
0 & 0 & 1
\end{bmatrix},
\qquad
f_{\text{logical}} = P f_{\text{stage}},
\qquad
q_{\text{stage}} = P^{\top} q_{\text{logical}}.
\]

The internal states remain in logical coordinates. Substituting
\(f_{\text{logical}} = Pu\) and \(y_{\text{stage}} = P^\top y\) gives
\[
B_{c,\text{stage}} = B_c P \in \R^{6\times 3},
\]
\[
C_{c,\text{stage}} = P^\top C_c \in \R^{3\times 6},
\]
while
\[
A_c(Y) \quad \text{remains unchanged.}
\]

\sourcenote{\texttt{kamtin-fp-model/03 Simulink gantry/main.m} lines 98--103
(\texttt{P} definition and \texttt{StageCoordinatesSystem = P.'*sys*P}).}

\end{derivationstep}

\begin{derivationstep}{ZOH discretization (Tóth complete method)}

Under Assumption~1 \cite{toth2010}, \(u(t)\) and \(Y(t)\approx Y[k]\) are held constant on
\[
[kT_d,(k+1)T_d).
\]

The continuous-time system therefore reduces to an LTI ODE on each interval. Solving exactly:
\[
x[k+1]
=
\underbrace{e^{A_c(Y[k])T_d}}_{A_d(Y[k])}x[k]
+
\underbrace{\left[\int_0^{T_d} e^{A_c(Y[k])\tau}\, d\tau\right] B_{c,\text{stage}}}_{B_d(Y[k])}u[k].
\]

The output equation is
\[
y[k] = C_{c,\text{stage}}x[k],
\qquad
D_d = 0.
\]

This result is a DT-LPV-SS model in the sense of Tóth \cite{toth2010} Definition~2, with
discrete scheduling signal $p_d(k) := p_c(kT_d) = Y(kT_d)$.

\medskip
\textbf{Point 1: self-scheduling property.}
$Y[k] = x_3[k]$ is not an external signal. It is the third component of the state vector
$x[k]$, updated by the recurrence itself. After each step the new state $x[k+1]$ contains
$x_3[k+1] = Y[k+1]$, which becomes the scheduling variable for the next interval. No
separate scheduling signal is needed: the model is self-scheduled through the state.

\medskip
\textbf{Point 2: chaining of frozen intervals.}
Tóth's complete method gives the exact solution for a single interval $[kT_d, (k+1)T_d)$
with $Y$ frozen at $Y[k]$. The full LPV simulation is formed by chaining these intervals:
\[
    x[k+1] = A_d(x_3[k])\,x[k] + B_d(x_3[k])\,u[k], \qquad k = 0,\,1,\,2,\,\ldots
\]
At each step $k$: (i) read $Y[k] = x_3[k]$ from the current state, (ii) compute
$A_d(Y[k])$ and $B_d(Y[k])$ via the augmented matrix exponential, (iii) advance the state.
The scheduling variable is updated automatically in step (iii) via the new $x_3[k+1]$.

\medskip
\textbf{Point 3: initial conditions.}
The model uses absolute coordinates (not deviations from an operating point). The initial
state is $x[0] = x_0$, where $x_3[0] = Y_0$ is the initial absolute payload position.

\medskip
\textbf{Point 4: backpropagation through time (training loop only).}
In the training loop the loss $\mathcal{L}$ must propagate gradients back through the
entire simulation sequence. The gradient chain at step $k$ is:
\[
    \frac{\partial \mathcal{L}}{\partial A_d(Y[k])} \cdot
    \frac{\partial A_d}{\partial Y[k]} \cdot
    \frac{\partial Y[k]}{\partial x[k]} \cdot
    \frac{\partial x[k]}{\partial x[k-1]} \cdots
\]
The middle term $\partial A_d / \partial Y[k]$ flows through
\texttt{torch.linalg.matrix\_exp}, which is differentiable. The full chain is therefore
traceable by autograd across all timesteps (BPTT). This requires that every operation in
Points~1--2 uses torch tensor arithmetic with no detach or numpy conversion.

\textbf{Quasi-LPV caveat.}
Tóth \cite{toth2010} notes this method is \emph{``only reasonable for the discretization of
LPV-SS representation with static dependence as dynamic dependence requires a higher-order
hold approach.''}
Our scheduling variable \(Y = x_3\) is a state (dynamic dependence), so a small residual
intra-sample error exists. This is accepted due to the 220:1 timescale separation between
\(Y\)-variation (\(\Delta Y \le 0.125\,\text{mm}\) per sample at \(f_s = 16\,\text{kHz}\)
PLTI rate \cite{etel_accuret}) and the plant dynamics (\(\sim\!1.6\,\text{ms}\) at 100~Hz bandwidth).
The maximum \(Y\)-velocity \(v_{\max} = 2\,\text{m/s}\) is taken from the ETEL Telica
specification \cite{etel_accuret} and confirmed in \texttt{main.m} (\texttt{vmax = 2}).

\textbf{Discrepancy: \texttt{main.m} vs.\ ETEL specification.}
The ETEL Telica specification \cite[p.~18]{etel_accuret} states three loop cycle times:
\begin{itemize}
    \item MLTI (manager loop): $400\,\mu\text{s}$ \quad ($f_\text{MLTI} = 2.5\,\text{kHz}$)
    \item PLTI (position loop): $50\,\mu\text{s}$ \quad ($f_\text{PLTI} = 20\,\text{kHz}$)
    \item CLTI (current loop): $50\,\mu\text{s}$ \quad ($f_\text{CLTI} = 20\,\text{kHz}$)
\end{itemize}
PLTI and CLTI run at the same rate and are synchronised: the force command computed by
the position loop is executed by the current loop at the same tick. There is no timing
cascade issue; $u[k]$ is held for exactly one PLTI period of $T_d = 50\,\mu\text{s}$.

\texttt{main.m} uses \texttt{fs = 16e3} ($T_d = 62.5\,\mu\text{s}$), which does
\textbf{not} match the ETEL specification ($T_d = 50\,\mu\text{s}$, $f_s = 20\,\text{kHz}$).
This discrepancy must be resolved before $T_d$ can be stated definitively. All numbers in
this document that depend on $T_d$ (in particular $\Delta Y_{\max}$ and the timescale
separation) use the \texttt{main.m} value until this is clarified. For reference, at
$f_s = 20\,\text{kHz}$:
\[
    \Delta Y_{\max} = v_{\max} \cdot T_d = 2 \times 50 \times 10^{-6} = 0.1\,\text{mm},
\]
which is smaller than the 0.125\,mm computed at 16\,kHz and only strengthens the ZOH
justification.

\sourcenote{\cite{toth2010} Assumptions 1 \& 2 (pages 5--6).
Loop rates: \cite[p.~18]{etel_accuret}.
Discrepancy flagged in \texttt{docs/decisions.md} D-012.}

\begin{tcolorbox}[
    colback=red!4, colframe=red!60!black, boxrule=0.9pt, arc=1.5mm,
    left=2mm, right=2mm, top=1mm, bottom=1mm,
    title={\bfseries Implementation requirement: scheduling wiring not yet done},
    coltitle=black, colbacktitle=red!15, fonttitle=\bfseries
]
The complete method produces $A_d(p_d(k))$ and $B_d(p_d(k))$ where
$p_d(k) = Y(kT_d) = x_3[k]$. This means at \textbf{every timestep} the training loop must:
\begin{enumerate}
    \item read $Y[k] = x_3[k]$ from the current state vector,
    \item pass it into \texttt{gantry\_lpv\_matrices\_torch(Y)} to obtain $A_d(Y[k])$
          and $B_d(Y[k])$,
    \item use those matrices for the state update $x[k+1] = A_d(Y[k])\,x[k] + B_d(Y[k])\,u[k]$.
\end{enumerate}
\textbf{This wiring does not exist yet.} The function
\texttt{gantry\_lpv\_matrices\_torch} is implemented and validated
(\texttt{scripts/gantry/gantry\_lpv\_torch.py}), but the caller that reads $x_3[k]$ from
the state and passes it in at each step belongs to \texttt{LPV\_Linear\_State\_Block.forward()},
which is not yet implemented (Step~3 of the project plan,
\texttt{docs/decisions.md} D-013).

\medskip
\textbf{Autograd compatibility.}
The proposed architecture (Points~1--4 above) is fully compatible with autograd.
Two gradient paths exist at each step $k$:
(a) direct: $\mathcal{L} \to x[k+1] \to A_d x[k] \to x[k]$, and
(b) self-scheduling: $\mathcal{L} \to x[k+1] \to A_d(Y[k]) \xrightarrow{\partial\,\mathrm{matrix\_exp}/\partial Y} Y[k] = x_k[2] \to x[k]$.
PyTorch accumulates gradients from both branches at $x[k]$ correctly.
Requirements: (i) extract scheduling variable as \texttt{Y\_k = x\_k[2]} (tensor, never \texttt{.item()}),
(ii) all state updates out-of-place, (iii) trajectory accumulated via \texttt{list} + \texttt{torch.stack}.
Gradient clipping (\texttt{clip\_grad\_norm\_}) is recommended due to the double gradient path.

\end{tcolorbox}

\end{derivationstep}

\begin{derivationstep}{Why the naive \(B_d\) formula fails}

When \(A_c\) is invertible, the integral simplifies as
\[
\int_0^{T_d} e^{A_c \tau}\, d\tau
=
A_c^{-1}(e^{A_c T_d}-I),
\]
which gives the convenient form \cite{toth2010} (eq.~9a):
\[
B_d = A_c^{-1}(A_d-I)\,B_{c,\text{stage}}.
\]

However, in our case \(A_c\) is singular, so \(A_c^{-1}\) does not exist and this formula is
undefined. Tóth acknowledges this in footnote~2: \emph{``Ac(p) is not required to be
invertible, but if it is, we can write the resulting DT description of the state-evolution
conveniently as (9a).''}

\sourcenote{\cite{toth2010} eq.~9a and footnote~2.
Design decision: \texttt{docs/decisions.md} D-015.}

\end{derivationstep}

\begin{derivationstep}{Augmented matrix exponential (general derivation)}

Form the augmented \((n+m)\times(n+m)=9\times 9\) matrix
\[
\mathcal{M}=
\begin{bmatrix}
A_c & B_{c,\text{stage}} \\
0_{3\times 6} & 0_{3\times 3}
\end{bmatrix}.
\]

For \(k\ge 1\),
\[
\mathcal{M}^k=
\begin{bmatrix}
A_c^k & A_c^{k-1}B_{c,\text{stage}} \\
0 & 0
\end{bmatrix}.
\]

Using the Taylor expansion,
\[
e^{\mathcal{M}T_d}
=
I+\sum_{k=1}^{\infty}\frac{\mathcal{M}^kT_d^k}{k!}
=
\begin{bmatrix}
\displaystyle I+\sum_{k=1}^{\infty}\frac{A_c^kT_d^k}{k!}
&
\displaystyle \sum_{k=1}^{\infty}\frac{A_c^{k-1}T_d^k}{k!}B_{c,\text{stage}}
\\[1.2em]
0 & I
\end{bmatrix}.
\]

The top-left block is
\[
e^{A_cT_d}=A_d.
\]

For the top-right block,
\[
\sum_{k=1}^{\infty}\frac{A_c^{k-1}T_d^k}{k!}
=
\int_0^{T_d}
\sum_{k=1}^{\infty}
\frac{A_c^{k-1}\tau^{k-1}}{(k-1)!}\, d\tau
=
\int_0^{T_d} e^{A_c\tau}\, d\tau.
\]

Therefore,
\[
\boxed{
e^{\mathcal{M}T_d}
=
\begin{bmatrix}
A_d & B_d \\
0 & I
\end{bmatrix}
}
\]
with no inversion of \(A_c\) required. Hence this method remains valid for singular \(A_c\).

This is exactly what \verb|torch.linalg.matrix_exp(M_aug * ts)| computes, and is identical
to what \verb|scipy.signal.cont2discrete(method='zoh')| uses internally.

\sourcenote{Van Loan (1978) for the augmented exponential identity.
Implementation: \texttt{scripts/gantry/gantry\_lpv\_torch.py} (torch, training loop) and
\texttt{scripts/gantry/gantry\_ss.py} (scipy, validation).
Design decision: \texttt{docs/decisions.md} D-015.}

\end{derivationstep}

\begin{thebibliography}{9}

\bibitem{toth2010}
R.~Tóth,
\textit{Modeling and Identification of Linear Parameter-Varying Systems},
Springer, 2010.
Relevant sections: Section~I (static vs.\ dynamic dependence, p.~2),
Assumptions~1 \& 2 (pp.~5--6), complete method eq.~9a and footnote~2.

\bibitem{etel_accuret}
ETEL~S.A.,
\textit{AccurET -- Operation \& Software, Version~V},
\texttt{literature/AccurET-Oper\&Soft-VerV.pdf}.
Page~18: PLTI = $50\,\mu\text{s}$ (20\,kHz), CLTI = $50\,\mu\text{s}$ (20\,kHz),
MLTI = $400\,\mu\text{s}$ (2.5\,kHz). Maximum Y-axis velocity \(v_{\max} = 2\,\text{m/s}\).
Note: \texttt{main.m} uses \texttt{fs = 16e3} (62.5\,$\mu$s) -- discrepancy with spec
must be resolved.

\end{thebibliography}

\end{document}
